\chapter{Постановка задачи}

Методы решения систем линейных алгебраических уравнений (СЛАУ) можно разделить на точные (прямые) методы, представляющие собой конечный алгоритм для получения решения, и итерационные методы, позволяющие получить решение с заданной точностью путем сходящихся бесконечных процессов. Вследствие неизбежных округлений результаты даже точных методов являются приближенными. При использовании итерационных методов добавляется погрешность самого метода.

Рассматривается система $n$ линейных алгебраических уравнений с $n$ неизвестными:
\begin{equation}
    \begin{cases}
        a_{11}x_1 + a_{12}x_2 + \ldots + a_{1n}x_n = b_1; \\
        a_{21}x_1 + a_{22}x_2 + \ldots + a_{2n}x_n = b_2; \\
        \dots\dots\dots\dots\dots\dots\dots\dots\dots\dots\dots \\
        a_{n1}x_1 + a_{n2}x_2 + \ldots + a_{nn}x_n = b_n,
    \end{cases}
    \label{eq:slau}
\end{equation}
или в матричной форме записи:
\begin{equation}
    Ax = b,
\end{equation}
где $A$~--- квадратная матрица коэффициентов системы порядка $n$; $b$~--- вектор свободных членов; $x$~--- искомый вектор неизвестных.

\section{Анализ задачи перед ее решением}

Перед непосредственным решением системы требуется провести анализ задачи:
\begin{enumerate}
    \item \textbf{Условие существования и единственности решения}:
    Система~\eqref{eq:slau} имеет единственное решение тогда и только тогда, когда определитель матрицы системы отличен от нуля:
    \begin{equation}
        \det(A) \neq 0.
    \end{equation}
    Если $\det(A) = 0$, то на определенном шаге вычислений все коэффициенты при неизвестных в строке обращаются в ноль ($0 \cdot x_1 + \ldots + 0 \cdot x_n$):
    \begin{itemize}
        \item если при этом соответствующий свободный член отличен от нуля ($b_i \neq 0$), получается противоречие $0 = b_i$, и система \textbf{не имеет решения};
        \item если соответствующий свободный член также равен нулю ($b_i = 0$), получается тождество $0 = 0$, и система имеет \textbf{бесконечное множество решений}.
    \end{itemize}
    \item \textbf{Условия применимости методов}:
    \begin{itemize}
        \item для метода исключения Гаусса: вычисления нельзя производить, если на каком-либо шаге диагональный элемент равен нулю. Кроме того, если элемент на главной диагонали мал, умножение на большие множители приводит к значительным ошибкам при вычитании. Данный недостаток устраняется выбором главного (наибольшего по модулю) элемента;
        \item для метода простой итерации: метод применим и сходится, если диагональные коэффициенты преобладают по абсолютной величине над остальными коэффициентами строки: $|a_{ii}| > \sum_{j \neq i} |a_{ij}|$. Если в исходной системе хотя бы некоторые диагональные коэффициенты не преобладают, систему с помощью линейного комбинирования уравнений заменяют эквивалентной системой с нужным свойством.
    \end{itemize}
\end{enumerate}

\section{Задачи работы}

В соответствии с индивидуальным заданием (вариант №~11) поставлены задачи:
\begin{enumerate}
    \item Решить систему линейных уравнений методом Гаусса с выбором главных элементов. Исключить аварийное завершение программы в случае деления на ноль, предусмотрев вывод характера решения (отсутствие решений или бесконечное множество решений);
    \item Методом простой итерации (для нечетного варианта) решить систему с точностью $\varepsilon = 10^{-5}$, предварительно приведя ее к виду, удобному для итераций. Для хранения неизвестных $x$ использовать два одномерных массива (на текущей и предыдущей итерациях) и проверять условие сходимости метода через нормы матриц, подчиненные октаэдрической и кубической нормам векторов.
\end{enumerate}
